\chapter{Getting help and sending feedback}\label{ch:help}

This chapter lists the places to look when you are stuck, explains the help
built into ProPrep, and shows how to report a problem or send a suggestion to
the developers. It ends with how to cite ProPrep in your work.

\section{Where to look first}\label{sec:help-first}

\begin{enumerate}
  \item \textbf{This manual.} Chapter~\ref{ch:troubleshooting} is organized by
  what you see on the screen. Chapter~\ref{ch:knownissues} lists the known
  limitations of version 1.23.0, each with a way around it. The Glossary
  (Appendix~\ref{app:glossary}) explains technical words, and the Menu map
  (Appendix~\ref{app:menumap}) shows where every tool is.
  \item \textbf{The screen itself.} ProPrep often tells you what is missing.
  A \texttt{⚠} line under a tool says what it needs, such as
  \texttt{⚠ Needs a topology — run Topology Generator first}. Many tool menus
  also show a \texttt{→ Suggestion:} line with the next step.
  \item \textbf{The \texttt{Using ...} line} at the top of each tool, which
  tells you which file the tool is working on
  (Section~\ref{sec:projectfolder-pickup}). Many surprises come from the wrong
  file.
  \item \textbf{The log files} in your project folder. When building a
  topology fails, the real reason is often only in \file{leap.log}
  (Chapter~\ref{ch:topology}).
\end{enumerate}

\section{Help inside ProPrep}\label{sec:help-inside}

\begin{itemize}
  \item \key{h} in the Full Menu shows a one-page summary of common workflows.
  In version 1.23.0 its menu codes are out of date; use the corrected list in
  Section~\ref{sec:menus-help}.
  \item In Workflow Mode, \key{u} then \key{h} shows a short explanation of the
  stages (Section~\ref{sec:menus-workflow}).
  \item Some tools have their own help item, such as \textsf{Display detailed
  help information} in the Force Field Parameterizer (\key{2c}).
\end{itemize}

\section{Seeing more detail: \texttt{--debug} and \texttt{--verbose}}\label{sec:help-debug}

When ProPrep stops with an error, it normally prints a short message such as:

\begin{screen}
Error: [Errno 2] No such file or directory: 'sub/deeper'
Run with --debug flag for more details
\end{screen}

To see more, start ProPrep again with one of these options and repeat what you
did:

\begin{tabularx}{\linewidth}{@{}lL@{}}
\toprule
\textbf{Option} & \textbf{What you get} \\
\midrule
\typed{--debug} & ProPrep runs normally, but some tools print extra
diagnostic lines, and an error shows the full technical report (a
\term{traceback}, the list of program steps that led to the error) instead of
the short message. You can switch the same mode on and off during a run with
\key{w} then \textsf{6.~Toggle debug mode}. \\
\typed{--verbose} or \typed{-v} & The screen fills with time-stamped technical
lines marked \texttt{DEBUG}, from start-up onwards. This is meant for the
developers. Use it only when they ask you to. \\
\bottomrule
\end{tabularx}

\begin{terminal}
proprep --full-menu --debug
\end{terminal}

\begin{tip}
When you report a problem, copy the error lines from the screen (from a
\typed{--debug} run if you can) into your report. In the browser version, a
plain-text copy of the whole session is also kept in
\file{~/.proprep/web_sessions/}.
\end{tip}

\section{Feedback \& Support (\texttt{f})}\label{sec:help-feedback}

The \key{f} item in the main menu opens a form that sends a bug report,
feature request, suggestion or question to the developers as an \term{issue}
on GitHub (a public web page where software projects collect reports). You can
send it anonymously.

\subsection*{The form, question by question}

\begin{steps}
  \item \prompt{Select feedback type (3)}: \key{1} Bug Report, \key{2} Feature
  Request, \key{3} Suggestion, \key{4} Question, \key{5} Other.
  \item \prompt{Module number (or type module name) (0)}: which tool it is
  about, or \key{0} for General. You can also type the tool's name. The list
  uses some older names (``PDB Loader'' is the Structure Loader) and does not
  include the newer tools, such as Docking or the Membrane Builder; choose
  \key{0} for those and name the tool in your text.
  \item \prompt{Email address ()}: optional. Leave it empty to stay anonymous,
  or type your email if you would like a reply.
  \item \prompt{Summary}: one line. Required.
  \item The detailed description: type one line at each \prompt{()} prompt;
  an empty line finishes. Required. For a problem, say what you were trying to
  do, what happened, and any error message.
  \item ProPrep shows a \textsf{Feedback Preview} and asks whether to submit:
\begin{screen}
═══ Feedback Preview ═══

  Type:                   Question
  Module:                 Structure Viewer
  Email:                  Anonymous
  Summary:                Manual test - do not submit

  Description:
  Testing the form only.


Submit this feedback to GitHub Issues? [y/n] (y): §ans«n»
Feedback not submitted.
\end{screen}
\end{steps}

\begin{watchout}
The default answer to \prompt{Submit this feedback to GitHub Issues? [y/n] (y)}
is \textbf{yes}. Pressing \key{Enter} sends your feedback. Type \key{n} if you
are not ready. To leave the form at any earlier point, press \key{Enter} at
\prompt{Summary}: ProPrep says \texttt{Summary is required.} and
\texttt{Feedback cancelled.}
\end{watchout}

\subsection*{Setting up GitHub (once)}

To post an issue, ProPrep needs a GitHub account and a \term{personal access
token}. This is a long code, like a password, that lets ProPrep post on your
behalf. The
form's opening box explains how to make one at
\url{https://github.com/settings/tokens} (``Generate new token (classic)'', with
only the \texttt{public\_repo} permission ticked).

Then store it in ProPrep:
\begin{steps}
  \item In the main menu, type \key{p} (Preferences), then \key{3}
  (\textsf{Configure GitHub feedback}). The form's opening box says
  ``2.~Configure GitHub''; in version 1.23.0 it is option 3.
  \item At \prompt{GitHub repository (owner/repo)}, press \key{Enter} to accept
  the suggested repository, or type another in the form \typed{owner/repo}.
  \item At \prompt{Token:}, paste the token and press \key{Enter}. What you type
  is hidden.
\end{steps}

ProPrep replies \texttt{✓ GitHub feedback configuration saved}. It does not
check the token with GitHub at this point; a wrong token shows up only when you
submit.

\begin{watchout}
The opening box says the token is stored securely, but in version 1.23.0 it is
saved as plain, readable text in \file{~/.proprep/settings.json}. Anyone who can
read your files can read it. Give the token only the \texttt{public\_repo}
permission, do not use this on a shared account, and delete the token on GitHub
when you no longer need it. \textsf{5.~Reset all preferences to defaults} under
\key{p} removes it from ProPrep.
\end{watchout}

\begin{behindscenes}
The repository the form suggests (\texttt{Mag14011/ProPrep}) is not the one
ProPrep's releases come from (\texttt{mjgplab/proprep}). If you are unsure where
a report should go, email the developers instead (next section).
\end{behindscenes}

\subsection*{Without a GitHub token}

If no token is set up and you answer \key{y}, ProPrep cannot post the issue.
According to its source, it then saves your feedback as a file in
\file{~/.proprep/feedback/}, which you can email to the developers. We did not
try this for the manual.

\section{Contacting the developers}\label{sec:help-contact}

ProPrep is developed by the Guberman-Pfeffer Lab at New Mexico State University.
The thank-you box shown when you exit ProPrep gives the contact address:

\begin{screen}
│ Feedback (positive or negative) is welcome at: mjgp@nmsu.edu           │
\end{screen}

A useful report contains:
\begin{itemize}
  \item the ProPrep version (\typed{proprep --version});
  \item your operating system (macOS, Linux or Windows with WSL);
  \item the command you started ProPrep with;
  \item what you did, what you expected, and what happened, with the exact
  error lines copied from the screen;
  \item if possible, the session file from your project folder
  (\file{proprep_session_<date>_<time>.json}), which lets the developers repeat
  your steps exactly.
\end{itemize}

\begin{tip}
A session file contains every answer you typed, including file names and
folder paths on your computer. Look through it before you send it if any of
that is private.
\end{tip}

\section{Citing ProPrep}\label{sec:help-cite}

If ProPrep helped your work, please cite it. The exit screen asks you to cite:

\begin{example}
Walker, A. \& Guberman-Pfeffer, M.J. (2026). ProPrep: An Interactive and
Instructional Interface for Proper Protein Preparation with AMBER. bioRxiv.
\url{https://doi.org/10.64898/2026.02.26.708365}
\end{example}

The ProPrep source code also comes with a citation file
(\file{CITATION.cff}), which asks you to cite both the software (ProPrep
version 1.23.0) and the accompanying article. Its entry for the article has a
newer title, ``ProPrep: (Metallo)protein MD Preparation That Keeps Users
Informed and in Control'', by Walker, Nielson, Cairns and Guberman-Pfeffer.
Check the ProPrep download page (\url{https://github.com/mjgplab/proprep}) for
the current reference before you publish.

ProPrep runs many other programs for you, among them AmberTools (tLEaP, sander,
ParmEd, cpptraj), PROPKA, packmol-memgen and AutoDock Vina. When your work uses
their results, cite them too, as their own documentation asks.
